% 编译: xelatex SL_gap_n1_research_summary.tex (两遍, -output-directory=build)
\documentclass[UTF8]{ctexart}
\usepackage{amsmath, amssymb, amsthm}
\usepackage[margin=2.5cm]{geometry}
\usepackage[hyphens]{url}
\usepackage[hidelinks]{hyperref}
\usepackage{booktabs}
\emergencystretch 3em

\newtheorem{theorem}{定理}[section]
\newtheorem{proposition}[theorem]{命题}
\newtheorem{lemma}[theorem]{引理}
\newtheorem{corollary}[theorem]{推论}
\newtheorem{remark}[theorem]{注}
\newtheorem{conjecture}[theorem]{猜想}

\title{相邻特征值间距 ($n=1$) 极端值: 严格证明研究总结 --- 部分结果}
\author{研究项目 (来源: BVE research, run R-20260805T000000Z-gapn1-a1b2c3)}
\date{2026-08-05}

\begin{document}
\maketitle

\begin{abstract}
本报告研究 Dirichlet 弦问题 $-y''=\lambda\rho y$ ($1\le\rho\le R$) 相邻特征值间距
$D:=\lambda_2-\lambda_1$ 的上下确界. 已严格证明: (i) \textbf{归约定理}:
上确界 (resp. 下确界) 在``单垒''(resp. ``单阱'') 两块密度族上达到;
(ii) \textbf{两块界}: 任意两块密度满足 $3\pi^2/R < \lambda_2-\lambda_1 < 3\pi^2$;
(iii) 对称族单次穿零在数值上等价于一个显式单变量不等式 (KEY LEMMA), 对称临界点唯一性
等价于三条分支引理 (O3a 引理 A--C). 本文如实记录已证部分、数值证据、失败路线与精确剩余缺口.
\end{abstract}

\tableofcontents

\section{问题设置与猜想}

\subsection{设定 (问题)}

Dirichlet 弦
\begin{equation}
	-y''(x) = \lambda\,\rho(x)\, y(x), \qquad y(0) = y(1) = 0,
	\qquad 1 \le \rho(x) \le R \text{ a.e.}
\end{equation}
$0 < \lambda_1(\rho) < \lambda_2(\rho)$ 为前两个特征值,
$D(\rho) := \lambda_2(\rho) - \lambda_1(\rho)$.

\medskip
\noindent\textbf{SUP}: 求 $S(R) := \sup\{D(\rho): 1\le\rho\le R\}$ 及全部极大化子.
\textbf{INF}: 求 $I(R) := \inf\{D(\rho): 1\le\rho\le R\}$ 及全部极小化子.

猜想 (基于数值, 本 run 目标): $S(R)$ 由对称三块垒 $[1,R,1]$ 达到:
$\rho = 1$ on $[0,u]\cup[1-u,1]$, $\rho = R$ on $(u,1-u)$, 其中 $u^*(R)$ 为
自洽方程 $f(u)=0$ 在 $(0,1/2)$ 内的唯一解; $I(R)$ 由对称三块阱 $[R,1,R]$ 达到.

\subsection{义务表}

\begin{center}
\begin{tabular}{lll}
\toprule
义务 & 内容 & 状态 \\
\midrule
O1 & 归约到两块族 & \textbf{PROVED} (见第 2 节) \\
O3b(1) & 两块界 $3\pi^2/R < D < 3\pi^2$ & \textbf{PROVED} (第 3 节) \\
O2 & 对称族单次穿零 & PARTIAL (缺 KEY LEMMA, 第 4 节) \\
O3a & 对称临界点唯一性 & PARTIAL (缺引理 A--C, 第 5 节) \\
O4/O5 & 对称临界值超越常数, 完全刻画 & OPEN \\
\bottomrule
\end{tabular}
\end{center}

\textbf{状态标签: RIGOROUS\_PARTIAL\_RESULT.} 归约与两块界已严格证明;
两个核心缺口 (KEY LEMMA, 引理 A--C) 均以显式闭式陈述, 数值全通过.
补上任意一个缺口即升级对应义务, 两个缺口齐备即得完整证明.

\section{归约定理 O1 (证明)}

\begin{theorem}[归约]
设 $\mathcal{K} = \{\rho: 1\le\rho\le R \text{ a.e.}\}$. 则
\begin{equation}
	S(R) = \max_{0\le a\le b\le 1} D\bigl(\rho = R \text{ on } (a,b),\; 1 \text{ 其余}\bigr),
	\qquad
	I(R) = \min_{0\le a\le b\le 1} D\bigl(\rho = 1 \text{ on } (a,b),\; R \text{ 其余}\bigr),
\end{equation}
且两块族上的极值均达到.
\end{theorem}

证明 (七步, 详见 run 内 \texttt{O1\_reduction\_draft.md}):
\begin{enumerate}
	\item \textbf{连续性}. Dirichlet Green 函数
		$G(x,t)=\min(x,t)(1-\max(x,t))$, $T_\rho f=\int_0^1 G(x,t)\rho(t)f(t)\,dt$
		为 $L^2$ 紧自伴算子, $\lambda_k(\rho)^{-1} = \mu_k(T_\rho)$;
		$\|T_\rho - T_\sigma\| \le \|G\|_\infty \|\rho-\sigma\|_2 \to 0$ 当
		$\|\rho-\sigma\|_1\to 0$ (因 $\|\rho-\sigma\|_2^2 \le 2R\|\rho-\sigma\|_1$),
		由 min-max 原理 $\lambda_k$ 关于 $L^1$ 连续.
	\item \textbf{紧性}. 两块族为紧集 (闭有界参数), 极值存在.
	\item \textbf{$N$ 跳参数化}. $\mathcal{K}_N$ 为
		$\{0\le x_1\le\cdots\le x_N\le 1\}\times[1,R]^{N+1}$ 的连续像.
	\item \textbf{移动跳点 (FH)}. 若 $\rho$ 在 $x_j$ 处由 $c_-$ 跳到 $c_+$,
		把跳点移到 $x_j+\varepsilon$, 由 Feynman--Hellmann 公式:
		\begin{equation}
			\frac{d}{d\varepsilon}D\big|_0 = (c_+ - c_-)\, f(x_j),
			\qquad f := \lambda_1 u_1^2 - \lambda_2 u_2^2
			\ (u_k \text{ 为 } L^2(\rho)\text{-归一化特征函数}).
		\end{equation}
	\item \textbf{$f$ 的结构} (AEH arXiv:2407.02459 引理 2.2 重导,
		与 $\rho$ 无关). Wronskian
		$W = u_1u_2' - u_1'u_2$ 满足 $W' = (\lambda_1-\lambda_2)\rho\,u_1u_2$,
		$W(0)=W(1)=0$, 故 $W<0$ on $(0,1)$, $v:=u_2/u_1$ 严格递减;
		因 $f=0 \iff v^2 = \lambda_1/\lambda_2$ 至多两解, $\{f>0\}$ 为包含
		$u_2$ 唯一零点 $z_0$ 的单区间.
	\item \textbf{极值点至多两跳}. 每个有效跳点 $x_j$ 满足 $f(x_j)=0$,
		由 $f$ 至多两个零点得跳数 $\le 2$.
	\item \textbf{稠密性与 bang-bang}. 阶梯函数在 $\mathcal{K}$ 中 $L^1$ 稠密;
		全局极大化子若在 $\{f>0\}$ 上取 $\rho<R$, 增大 $\rho$ 会严格增大 $D$,
		矛盾, 故 $\rho = R$ a.e. on $\{f>0\}$, $\rho=1$ a.e. on $\{f<0\}$;
		由第 5 步, $\{f>0\}$ 为单区间, 极大化子即单垒两块配置. INF 同理 (符号翻转).
\end{enumerate}

注: 该归约给出``任意 $\rho$ 满足
$D(\rho) \le \max_{\text{单垒}}$ (resp. $\ge \min_{\text{单阱}}$)''的严格不等式链,
不依赖对称性先验假设.

\section{两块界定理 O3b(1) (证明)}

\begin{theorem}[两块界]
对两块密度 $\rho = 1$ on $[0,t]$, $R$ on $(t,1]$ (重右),
\begin{equation}
	\frac{3\pi^2}{R} \;<\; D(t) \;<\; 3\pi^2 \qquad \forall\, t\in(0,1),\ R>1.
\end{equation}
\end{theorem}

证明 (Agent C, 详见 run 内 \texttt{agentC\_O3b\_boundary.md}; 重左由镜像 $x\mapsto1-x$ 归入):
\begin{enumerate}
	\item \textbf{相位坐标}. 令 $\mu=\sqrt R$, $c = \mu(1-t)/t$,
		$\theta(x)=\arctan(\mu\tan x)$ 连续分支 ($\theta(x+\pi)=\theta(x)+\pi$),
		$x_1<x_2$ 为 $\theta(x) + cx = k\pi$ ($k=1,2$) 的前两根, 则
		$\lambda_k = x_k^2(\mu+c)^2/\mu^2$, 于是只需证
		$3\pi^2 < W(\mu,c) := (\mu+c)^2(x_2^2-x_1^2) < 3\pi^2\mu^2$.
	\item \textbf{下界}: 由 $\theta' < \mu$ 严格, 从 secular 方程得
		$x_1 > \pi/(\mu+c)$, $x_2 - x_1 > \pi/(\mu+c)$, 故
		$W > (\mu+c)^2\cdot\frac{\pi}{\mu+c}\cdot\frac{3\pi}{\mu+c} = 3\pi^2$.
	\item \textbf{上界}, 分三区:
		\begin{itemize}
			\item $c \ge 1$: $\theta(x)\ge x$ on $(0,\pi/2)$ 得
				$x_1 \le \pi/(1+c)$; 由 $\theta'\ge 1/\mu$ 得
				$x_2 - x_1 \le \pi\mu/(1+\mu c)$, $x_2 \le 2\pi\mu/(1+\mu c)$;
				得到 $W \le 3\pi^2\mu^2 G(\mu,c)$,
				$G(1,c)=1$ 且 $dG/d\mu < 0$ (sympy 精确分解:
				$dG/d\mu$ 分子为 $-(s+t+2)P(s,t)$, $s=\mu-1$, $t=c-1$,
				$P$ 为全正系数多项式且 $\ge 4$).
			\item $1/3 \le c \le 1$: 弦/凸性论证给出
				$x_2^2-x_1^2 \le 3\pi^2/(1+c)^2$, 故
				$W < 3\pi^2\mu^2$ (因 $(\mu+c)^2 < \mu^2(1+c)^2$ for $\mu>1$).
			\item $0 < c \le 1/3$: 令 $\varepsilon_k = k\pi - x_k$,
				$s_k = \tan(cx_k) = \mu\tan\varepsilon_k$, 证明 $W' < 0$ 于该区
				(用 $h(s,x)=s^2x^2/(\mu(\mu+c)+s^2(1+\mu c))$ 型估计), 故
				$W(\mu,c) < W(\mu,0) = 3\pi^2\mu^2$.
		\end{itemize}
\end{enumerate}

验证: $R \in \{1.05,\dots,10^4\}$ 上 4000 点网格零违例
(下界余量 $+1.25\times10^{-8}$, 上界余量 $+1.28\times10^{-6}$);
相位恒等式精度 $10^{-13}$; $W'<0$ 用 mpmath 60 位确认.

\section{O2: 对称族单次穿零 --- KEY LEMMA}

\subsection{设定}

对称垒 $[1,R,1]$ (参数 $u \in (0,1/2)$), 半问题 (在 $1/2$ 处奇/偶):
偶 (Neumann at $1/2$) secular 方程
$\tan(s_1u)\tan(s_1 q v) = 1/q$, 奇 (Dirichlet at $1/2$):
$q\tan(s_2u) + \tan(s_2qv) = 0$.
(更正: 任务初版的奇 secular 方程 $\tan(s_2u)\tan(s_2qv)=-q$ 不成立;
正确形式为 $q\tan(s_2u)+\tan(s_2qv)=0$, 已对转移矩阵验证 $10^{-12}$.)

记号: $q=\sqrt R$, $v=1/2-u$, $c = qv/u$, $u = q/(2(c+q))$, $\alpha_k = s_ku$,
$\beta_k = s_k q v = c\alpha_k$. 由 Feynman--Hellmann:
\begin{equation}
	\frac{dD}{du} = -2(R-1)\, f_{\rm sym}(u),
	\qquad f_{\rm sym}(u) = \lambda_1 u_1(u)^2 - \lambda_2 u_2(u)^2 .
\end{equation}
定义
\begin{equation}
	M(\alpha;c) := \frac{q(q^2-1)\,\alpha^2\sin^2\alpha}{q + c\,\Phi(\alpha)},
	\qquad \Phi(\alpha) := \cos^2\alpha + q^2\sin^2\alpha,
\end{equation}
\begin{equation}
	M_k(c) := M(\alpha_k(c);c), \qquad F(c) := M_1 - M_2,
\end{equation}
其中 $\alpha_1(c)$ ($\in(0,\pi/2)$) 与 $\alpha_2(c)$ ($\in(0,\pi)$) 为
$\beta = c\alpha$ 与偶/奇 secular 曲线的交点. 数值验证 (Agent A, 精度
$10^{-12}$):
\begin{equation}
	f_{\rm sym}(u) = \frac{2(c+q)}{q\,u^2(q^2-1)}\,F(c), \qquad
	D'(c) = \frac{8}{q^2}(c+q)\,F(c).
\end{equation}
故 $f_{\rm sym}$ 与 $F$ 同号, $D_{\rm sym}$ 的极值点与 $F$ 的零点一一对应.

\subsection{已证命题}

\begin{theorem}[O2 结构, 部分]
(i) 若 $c\ge 1$: $F(c) < 0$ (引理 1: $\phi_c(\alpha)=\alpha^2\sin^2\alpha/(q+c\Phi(\alpha))$
在 $(0,\pi/2)$ 严格递减, 用 $c/(q+c\Phi) < 1/\Phi$ 与
$(q^2-1)t^2 \le 1+q^2t^2$).
(ii) 若 $c\in[1/2,1]$: $F(c) < 0$ (用 $c=1/2$ 处
$\alpha_1 = \pi-\alpha_2 =: \alpha_0$, $\sin(\alpha_0/2)=1/\sqrt{2(q+1)}$).
(iii) 端点数据: $F(0+) = (q^2-1)\pi^2/4 > 0$; $F(1) < 0$;
$f_{\rm sym}(1/2) = 2\pi^2 > 0$; $f_{\rm sym}(u) \sim -30\pi^4u^2/R^2$ ($u\to0^+$).
(iv) 若 KEY LEMMA 成立, 则 $f_{\rm sym}$ 在 $(0,1/2)$ 有唯一零点
$u^*$, 符号由 $-$ 变 $+$, $D_{\rm sym}$ 先增后减, $u^*$ 为对称族
唯一极大点.
\end{theorem}

\subsection{KEY LEMMA (未证, 核心缺口)}

\begin{conjecture}[KEY LEMMA]
对所有 $q>1$, $c\in(0,1/2)$:
\begin{equation}
	\frac{d}{dc}\log\frac{M_1}{M_2} < 0.
\end{equation}
等价形式: $G(\alpha_1(c);c) < G(\alpha_2(c);c)$, 其中
\begin{equation}
	G(\alpha;c) := -\frac{\Phi(\alpha)W(\alpha)}{q+c\Phi(\alpha)}
		+ \frac{2c\,\alpha\,\Phi(\alpha)(q^2-1)\sin\alpha\cos\alpha}{(q+c\Phi(\alpha))^2},
	\qquad W(\alpha) := 3 + 2\alpha\cot\alpha .
\end{equation}
\end{conjecture}

\textbf{本次推进 (分解)}. 令 $\alpha_2 = \pi-\gamma$, 则 $G_2-G_1$ 可分解为
\begin{equation}
	G_2 - G_1 = (A-C) + (B-D),
\end{equation}
\begin{equation}
	A-C = \frac{\Phi_1 W_1}{q+c\Phi_1} - \frac{2c(q^2-1)\alpha_1\Phi_1\sin\alpha_1\cos\alpha_1}{(q+c\Phi_1)^2},
	\qquad
	B-D = -\frac{\Phi_2 W_2}{q+c\Phi_2}
		+ \frac{2c(q^2-1)\alpha_2\Phi_2\,|\sin\alpha_2\cos\alpha_2|}{(q+c\Phi_2)^2}.
\end{equation}
($A-C$ 用偶侧, $B-D$ 用 $\alpha_2=\pi-\gamma$ 与 $|\sin\alpha_2\cos\alpha_2|$.)
数值 (全网格 $q\in[1.00025, 10^6]$, $c\in(0,1/2)$):
\begin{center}
\begin{tabular}{lcc}
\toprule
量 & 最小值 & 达于 \\
\midrule
$A-C$ & $2.80613$ & $q\to1^+$, $c\to1/2^-$ (精确角点极限) \\
$B-D$ & $-0.38773$ & $q\to1^+$, $c\to1/2^-$ (精确角点极限) \\
$G_2-G_1 = (A-C)+(B-D)$ & $2.41840$ & $q\to1^+$, $c\to1/2^-$ (即 $q=1$ 基值) \\
\bottomrule
\end{tabular}
\end{center}
$q=1$ (即 $R=1$) 基值闭式 ($\alpha_1=\pi/(2(1+c))$,
$\alpha_2=\pi/(1+c)$, $\Phi\equiv1$, $(q^2-1)=0$):
\begin{equation}
	(A-C)\big|_{q=1} = \frac{W(\alpha_1)}{1+c},
	\qquad
	(B-D)\big|_{q=1} = -\frac{W(\alpha_2)}{1+c},
\end{equation}
由 $W'(\alpha) = 2(\sin\alpha\cos\alpha - \alpha)/\sin^2\alpha < 0$
得 $W(\alpha_2) < W(\alpha_1)$, 故
\begin{equation}
	(G_2-G_1)\big|_{q=1} \ge \frac{W(\pi/3) - W(2\pi/3)}{3/2}
		= \frac{4\pi}{3\sqrt3} = 2.41840\ldots > 0 .
\end{equation}
角点极限 (精确): 在 $q\to1^+$, $c\to1/2^-$ 处 $\alpha_1\to\pi/3$, $\alpha_2\to2\pi/3$,
$A-C \to W(\pi/3)/(3/2) = 2.80613\ldots$, $B-D \to -W(2\pi/3)/(3/2) = -0.38773\ldots$,
和 $\to 4\pi/(3\sqrt3) = 2.41840\ldots$.

\textbf{重要更正 (独立复核 2026-08-05): 逐项 $q$-单调性路线被否证.} 细网格复核显示 $A-C$ 对 $q$ 单调递增 (全采样通过), 但 $B-D$ 	extbf{不}对 $q$ 单调:
反例 $c=0.01$, $q\colon 5000\to20000$ 时 $B-D\colon 199.79\to193.99$ (递减); $c\le0.1$ 时均递减, $c\ge0.3$ 时才递增. 因此交接稿中``$d(B-D)/dq\ge0$ 全网格成立''的陈述不成立, ``分解 + 逐项 $q$-单调''的闭环方案作废. 但和的整体下界数值稳健: 全网格 $q\in[1.00025,10^6]$, $c\in(0,1/2)$ 扫描
$G_2-G_1 \ge 2.41840$, 最小值恰在角点 $q\to1^+$, $c\to1/2^-$ (即 $q=1$ 基值). KEY LEMMA 仍开放; 新的可能机制: $A-C$ 的 $q$-单调性 (可证方向) 与 $B-D$ 在
$A-C$ 增量区域的大正互补性, 或 $G_2-G_1$ 整体的解析下界.

\textbf{否证清单} (Agent A): (1) $F$ 在 $(0,1)$ 单调/凸 -- 假;
(2) $M_1/M_2$ 全局递减 -- 假 ($R\ge10$ 时 $(1/2,1)$ 内有非单调块);
(3) $G(\alpha;c)$ 对 $\alpha$ 单调 -- 假; (4) 符号二分 $G_1<0<G_2$ --
假 ($R<4$ 时 $c\to1/2$ 处 $G_2<0$); (5) $D$ 凹凸/单峰性 -- 假;
(6) Wronskian 迁移论证 -- 不适用. 补充: (7) $G_2-G_1$ 对 $c$ 单调 -- 假
($R=10^4$ 时 $c\approx0.44$ 处 $d(G_2-G_1)/dc > 0$);
(8) 逐项 $q$-单调性 $d(B-D)/dq\ge0$ -- 假 ($c\le0.1$ 时 $B-D$ 随 $q$ 递减),
故``分解 + 逐项 $q$-单调''闭环方案作废 (见 4.3 更正).

\section{O3a: 对称临界点唯一性}

\subsection{设定}

单垒 $\rho_{(a,b)} = R$ on $(a,b)$, 1 其余. 记
$R_1(a,b):=f(a;a,b)$, $R_2(a,b):=f(b;a,b)$. 符号一致临界点指
$f(a)=f(b)=0$ 且符号模式 $-,+,-$ 于 $(a,b)$; 反射 $\sigma(a,b)=(1-b,1-a)$.

\subsection{已证定理 (Agent B, 恒等式 T3 独立复核 \texorpdfstring{$2.7\times10^{-6}$}{2.7e-6})}

\begin{theorem}[O3a 结构]
(i) 符号一致临界点 $\iff$ $T(a,b):=(f\text{ 的两零点})$ 的不动点
$\iff$ $(R_1,R_2)=0$ 且 $(a,b)$ 为``好根''.
(ii) $T\circ\sigma = \sigma\circ T$; 不动点唯一 $\implies b=1-a$.
(iii) 精确恒等式 $dR_1/db = -dR_2/da$ (由 FH 公式与 Schwarz 定理,
$D$ 在 $(a,b)$ 上 $C^2$).
(iv) 条件唯一性: 若存在区间上两条 $C^1$ 好分支
$b=g_1(a)$, $b=g_2(a)$ 穿过一切不动点, 且 $g_1' > g_2'$,
$(g_1-g_2)$ 两端异号, 则 $T$ 至多一个不动点 (进而唯一且对称).
\end{theorem}

数值 (R=4): 唯一好根 $(0.451485465757, 0.548514534243)$, 对称;
$J_T$ 谱半径 $0.5611 < 1$ (局部稳定). 但 $R=100$ 时
$\rho(J_T)=1.642$ 且 $T$ 有真 2-周期, 故 $T$ 非全局压缩.
对 13 个 $R$ 值 ($1.02$ 至 $1000$) 的 126 网格点扫描无例外; 大 $R$ 需排除``伪根'' ($R=50/100$ 处 $(0.002,0.997)$ 型残差零点
非符号一致, 见 Agent B 3.5).

\section{失败路线}

\begin{description}
	\item[对称形状最优 (Route B) 失败] 比较函数 $g(x)=4\sin^2(2\pi x)-\sin^2(\pi x)$
		在 $x\approx0.2400, 0.7600$ 处穿零, 非单区间. 窄垒 ($w\le0.2$)
		时 $D(m,w)$ 的 $m$ 最优确在 $m=1/2$, 但 $w\ge0.3$ 失败.
		教训: 逐点比较函数不满足所需符号结构; 需全局变分论证.
	\item[$\tilde\rho$ 单调路线 (Route F) 失败] 对
		$\rho_k(u)=\lambda_2u_2^2/(\lambda_1u_1^2)$ 的单调性:
		$R=2$ 通过, $R=4$ 有 20/99 反例, $R=10$ 有 42/99. 失效.
	\item[压缩论证] $T$ 非全局压缩 ($R=100$ 谱半径 1.64, 有真 2-周期);
		``欧氏度量 2.3''型论证失败 (最大比 1.59).
	\item[对 $c$ 单调路线] $G_2-G_1$ 对 $c$ 非单调 (见 4.3).
	\item[逐项 $q$-单调路线 (B-D) 失败] $A-C$ 对 $q$ 单调递增 (细网格验证),
		但 $B-D$ 不单调: $c=0.01$, $q\colon5000\to20000$ 时 $B-D\colon199.79\to193.99$.
		交接稿中``$d(B-D)/dq\ge0$''陈述不成立, 该闭环方案作废 (见 4.3 更正);
		和的整体下界 $G_2-G_1\ge2.4184$ 数值稳健但证明仍开放.
	\item[文献] Wiley/ScienceDirect 403; zbMATH 全文受限; 未发现直接覆盖
		Dirichlet 盒类 $1\le\rho\le R$ 间距极值的已发表结果 (文献审计 Agent Pasteur).
	\item[数值陷阱] 初版奇 secular 方程 (tan 形式) 错误;
		大 $R$ 下粗网格根误报 (Agent C).
\end{description}

\section{数值结果}

\begin{center}
\begin{tabular}{c|cc|cc}
\toprule
$R$ & SUP $u^*$ & SUP $D$ & INF $u^*$ & INF $D$ \\
\midrule
1.05 & 0.420835 & 29.7107 & 0.418296 & 28.1025 \\
1.1  & 0.422035 & 29.8085 & 0.417076 & 26.7380 \\
1.5  & 0.429832 & 30.4730 & 0.408814 & 19.1954 \\
2    & 0.436696 & 31.1023 & 0.401037 & 14.1278 \\
3    & 0.445665 & 31.9922 & 0.390127 & 9.1897 \\
4    & 0.451485 & 32.6140 & 0.382598 & 6.7845 \\
10   & 0.466931 & 34.4513 & 0.361313 & 2.6089 \\
100  & 0.488529 & 37.5470 & 0.334804 & 0.2509 \\
1000 & 0.496261 & 38.8254 & 0.330446 & 0.0250 \\
\bottomrule
\end{tabular}
\end{center}

数值方法 (见 \texttt{scripts/num\_gapn1\_landscape.py} 等): 对
$R\in\{1.05,\dots,1000\}$ 每族 SUP/INF 唯一临界点 = 对称点, 无例外.
极限: $R\to1^+$: $u^*\to\arccos(1/4)/\pi = 0.419569$; $R\to\infty$:
SUP $u^*\to1/2$, $D\to4\pi^2$; INF $u^*\to0.32992$, $D\cdot R\to24.943866$.
$dD/du = -2(R-1)f_{\rm sym}$ 与 FH 验证 $10^{-6}$--$10^{-13}$.

\section{后续方向}

\begin{enumerate}
	\item \textbf{KEY LEMMA}: 逐项 $q$-单调闭环已否证 ($B-D$ 非单调),
		新方向: (a) 证明 $A-C$ 的 $q$-单调性 (数值全通过);
		(b) 在 $B-D$ 递减区域 ($c\le0.1$) 利用 $B-D$ 的大正性与 $A-C$ 增量互补,
		或直接对和 $G_2-G_1$ 建立解析下界; (c) 角点极限已精确化 ($q\to1^+$,
		$c\to1/2^-$ 处 $G_2-G_1\to4\pi/(3\sqrt3)$), 可作展开基点.
		当前最高杠杆缺口; 建议符号计算 (sympy) 或解析单调性.
	\item \textbf{O3a 引理 A}: $g_1'-g_2'$ 的 $R$-一致正下界
		(数值从 $42.78$ (R=1.05) 降到 $0.287$ (R=100), 需 R 无关下界).
	\item \textbf{INF 极限}: 严格证明 $R\to\infty$ 时 $u^*\to u_\infty$
		且 $D\cdot R\to (a^2-\pi^2/4)/u_\infty^2 = 24.943866$.
	\item \textbf{合流}: O1+O2+O3a+O3b 齐备后撰写
		\texttt{SL\_gap\_n1\_proof.tex} 完整证明.
	\item \textbf{推广}: Neumann/混合边界, $q\ge0$ 势, $p$-Laplacian.
\end{enumerate}

\section{涉及到的数学知识}

\begin{description}
	\item[Sturm--Liouville 谱理论] 特征值简单性, 特征函数振荡定理
		(第 $k$ 个特征函数恰有 $k-1$ 个内零点).
	\item[Feynman--Hellmann 公式] 参数族 $\rho_\varepsilon$ 下
		$d\lambda_k/d\varepsilon = -\lambda_k\int (d\rho/d\varepsilon)u_k^2\,dx$.
	\item[Wronskian 与比值单调性] $v = u_2/u_1$ 严格递减
		$\implies$ $f$ 至多两零点, 单正区间.
	\item[min-max 原理与 Green 算子] 特征值倒数 $\mu_k(T_\rho)$,
		紧自伴算子扰动连续性.
	\item[Pr\"ufer 相位 / secular 方程] 两块问题化为
		$\theta(x)=\arctan(\mu\tan x)$ 相位方程.
	\item[bang-bang 变分] 逐点 FH 导数符号 $\implies$ 极值密度取端点值.
	\item[Schwarz 定理与闭 1-形式] 混合偏导交换
		$\implies$ 残差恒等式 $dR_1/db=-dR_2/da$.
	\item[归一化恒等式] 交界面处
		$u_k(u,u)^2=\tan^2\alpha_k/(1/2+w\tan^2\alpha_k)$.
	\item[分解论证] $G_2-G_1$ 拆为 $(A-C)+(B-D)$,
		逐项 $q$-单调 + $q=1$ 基值闭式.
\end{description}

\begin{thebibliography}{9}
\bibitem{aeh} A. Ahrami, M. El Allali, E. M. Harrell II, \emph{The minimal gap
	of the vibrating string with monotone density}, arXiv:2407.02459.
	\url{https://arxiv.org/abs/2407.02459}
\bibitem{keller} J. B. Keller, \emph{The minimum ratio of two eigenvalues},
	SIAM J. Appl. Math. 31 (1976) 485--491.
	\url{https://doi.org/10.1137/0131042}
\bibitem{mw} T. Mahar, B. Willner, \emph{An extremal eigenvalue problem},
	Comm. Pure Appl. Math. 29 (1976) 517--529.
	\url{https://doi.org/10.1002/cpa.3160290505}
\bibitem{wm} B. Willner, T. Mahar, \emph{Extremal eigenvalue problems for
	Sturm--Liouville problems}, SIAM J. Math. Anal. 13 (1982) 984--992.
\bibitem{ckll} C.-H. Cheng, K.-Y. Kung, C.-K. Law, W.-C. Lian, \emph{The
	dual eigenvalue problems and the gaps of the vibrating string},
	Comput. Math. Appl. 60 (2010) 2135--2143.
\bibitem{ab} M. S. Ashbaugh, R. D. Benguria, \emph{Optimal lower bound for the
	gap between the first two eigenvalues of one-dimensional Schr\"odinger
	operators with symmetric single-well potentials}, Proc. AMS 105 (1989)
	419--424.
\bibitem{he} S. Hedhly, \emph{Sharp upper bound for the ratio of the first two
	eigenvalues of the vibrating string with single-well density},
	arXiv:2111.01728. \url{https://arxiv.org/abs/2111.01728}
\bibitem{sun} H. Sun, \emph{On the minimum eigenvalue gap for the vibrating
	string}, J. Math. Anal. Appl. 516 (2022) 126513.
\bibitem{qi} C. Qi, H. Li, H. Xie, \emph{Extremal eigenvalue gaps of the
	vibrating string with symmetric densities}, Qual. Theory Dyn. Syst. 19
	(2020) 12.
\bibitem{hor} M. Horvath, \emph{On the first two eigenvalues of a Sturm--
	Liouville operator}, Proc. AMS 131 (2003) 799--803.
\end{thebibliography}

\end{document}